\label{sec:exp_setup}
\subsection{Research Questions}

In this paper, we mainly investigate the following research questions through experiments.

\begin{itemize}

    \item \textbf{RQ1}: How effective is MIST in improving code coverage for DBMSs compared to baseline approaches?

    \item \textbf{RQ2}: How does MIST perform across different DBMS modules (i.e., \textit{Parser}, \textit{Optimizer}, \textit{Executor}, \textit{Storage})?
    
    \item \textbf{RQ3}: What are the impacts of the two stages in MIST on code coverage improvement?
       
\end{itemize}

\subsection{Studied DBMSs}

Following prior studies~\cite{zhong2025scaling,zhong2025testing}, we evaluate MIST on three widely-used open-source DBMS representing diverse architectural designs and SQL dialect characteristics.

\begin{itemize}

\item \textbf{DuckDB~\cite{raasveldt2019duckdb}:} An in-process analytical database designed for OLAP workloads with columnar storage and vectorized execution. It supports modern SQL features including window functions and CTEs.

\item \textbf{PostgreSQL~\cite{stonebraker1986design}:} A mature relational database system with a sophisticated query optimizer and advanced features such as materialized views and JSON operations.

\item \textbf{SQLite~\cite{owens2010sqlite}:} A lightweight embedded database engine widely deployed in mobile and embedded systems. It implements a comprehensive SQL dialect with unique features such as virtual tables.

\end{itemize}

These three systems collectively cover a broad spectrum of DBMS architectures (in-process analytical, client-server relational, and embedded transactional, respectively), enabling us to comprehensively evaluate MIST's generalizability across different SQL dialects and execution models.

% \subsection{Studied LLMs and }

% To evaluate MIST's effectiveness across different model capacities, we select four open-source LLMs spanning different parameter scales and model families: Qwen2.5-7B, Llama3.1-8B, Qwen2.5-14B, and Qwen2.5-32B. These models represent lightweight to medium-sized LLMs that are practical for local deployment in industrial environments, aligning with the resource constraints commonly faced in real-world DBMS testing scenarios. 

\subsection{Studied LLMs and Baseline}

To evaluate MIST's effectiveness across different model capacities, we select four open-source LLMs spanning different parameter scales and model families: Qwen2.5-7B, Llama3.1-8B, Qwen2.5-14B, and Qwen2.5-32B. These models represent lightweight to medium-sized LLMs that are practical for local deployment in industrial environments, aligning with the resource constraints commonly faced in real-world DBMS testing scenarios.

For baseline comparison, we select Fuzz4All~\cite{fuzz4all} as our primary baseline. To the best of our knowledge, there is currently no existing work that uses LLMs as complete test case generators for DBMS testing. Inspired by ShQveL~\cite{zhong2025testing}, we adopt the state-of-the-art LLM-based framework Fuzz4All for fair comparison. Fuzz4All is the first universal fuzzer that leverages LLMs to generate test inputs across multiple programming languages and systems.  

\subsection{Metrics}

We evaluate MIST using two categories of coverage metrics to comprehensively assess its effectiveness in exercising DBMS code.

% Please add the following required packages to your document preamble:
% \usepackage{multirow}
% \usepackage[table,xcdraw]{xcolor}
% Beamer presentation requires \usepackage{colortbl} instead of \usepackage[table,xcdraw]{xcolor}
\begin{table*}[]
\caption{Overall code coverage comparison between baseline and MIST.}
\label{table:rq1}
\begin{tabular}{c|c|ccc|ccc|ccc}
\toprule
 &  & \multicolumn{3}{c|}{\textbf{DuckDB}} & \multicolumn{3}{c|}{\textbf{PostgreSQL}} & \multicolumn{3}{c}{\textbf{SQLite}} \\ \cline{3-11} 
\multirow{-2}{*}{\textbf{Model}} & \multirow{-2}{*}{\textbf{Approach}} & \textbf{Line} & \textbf{Function} & \textbf{Branch} & \textbf{Line} & \textbf{Function} & \textbf{Branch} & \textbf{Line} & \textbf{Function} & \textbf{Branch} \\ \midrule
 & Fuzz4All & 29.9\% & 29.7\%  & 17.0\% & 23.1\% & 25.2\% & 16.2\% & 29.9\% & 38.9\%  & 22.5\% \\
\multirow{-2}{*}{Qwen2.5-7B} & \cellcolor[HTML]{ECF4FF}\textbf{MIST} & \cellcolor[HTML]{ECF4FF}\textbf{42.0\%} & \cellcolor[HTML]{ECF4FF}\textbf{38.2\%} & \cellcolor[HTML]{ECF4FF}\textbf{24.5\%} & \cellcolor[HTML]{ECF4FF}\textbf{30.3\%} & \cellcolor[HTML]{ECF4FF}\textbf{38.6\%} & \cellcolor[HTML]{ECF4FF}\textbf{21.3\%} & \cellcolor[HTML]{ECF4FF}\textbf{39.7\%} & \cellcolor[HTML]{ECF4FF}\textbf{46.1\%} & \cellcolor[HTML]{ECF4FF}\textbf{31.3\%} \\ \midrule
 & Fuzz4All & 32.3\% & 31.7\% & 18.3\% & 24.2\% & 31.8\% & 16.9\% & 28.3\% & 36.6\% & 21.1\% \\
\multirow{-2}{*}{Llama3.1-8B} & \cellcolor[HTML]{ECF4FF}\textbf{MIST} & \cellcolor[HTML]{ECF4FF}\textbf{43.2\%} & \cellcolor[HTML]{ECF4FF}\textbf{39.3\%} & \cellcolor[HTML]{ECF4FF}\textbf{25.4\%} & \cellcolor[HTML]{ECF4FF}\textbf{29.3\%} & \cellcolor[HTML]{ECF4FF}\textbf{37.7\%} & \cellcolor[HTML]{ECF4FF}\textbf{20.5\%} & \cellcolor[HTML]{ECF4FF}\textbf{39.9\%} & \cellcolor[HTML]{ECF4FF}\textbf{45.2\%} & \cellcolor[HTML]{ECF4FF}\textbf{32.1\%} \\ \midrule
 & Fuzz4All & 26.4\% & 26.8\% & 14.6\% & 23.6\% & 31.3\% & 16.4\% & 21.6\% & 30.3\% & 15.6\% \\
\multirow{-2}{*}{Qwen2.5-14B} & \cellcolor[HTML]{ECF4FF}\textbf{MIST} & \cellcolor[HTML]{ECF4FF}\textbf{42.2\%} & \cellcolor[HTML]{ECF4FF}\textbf{38.6\%} & \cellcolor[HTML]{ECF4FF}\textbf{24.7\%} & \cellcolor[HTML]{ECF4FF}\textbf{31.4\%} & \cellcolor[HTML]{ECF4FF}\textbf{39.9\%} & \cellcolor[HTML]{ECF4FF}\textbf{22.2\%} & \cellcolor[HTML]{ECF4FF}\textbf{38.1\%} & \cellcolor[HTML]{ECF4FF}\textbf{44.5\%} & \cellcolor[HTML]{ECF4FF}\textbf{30.1\%} \\ \midrule
 & Fuzz4All & 25.8\% & 28.0\% & 18.6\% & 20.9\% & 28.4\% & 14.4\% &26.5\%  &34.7\%  & 19.8\% \\
\multirow{-2}{*}{Qwen2.5-32B} & \cellcolor[HTML]{ECF4FF}\textbf{MIST} & \cellcolor[HTML]{ECF4FF}\textbf{41.8\%} & \cellcolor[HTML]{ECF4FF}\textbf{38.6\%} & \cellcolor[HTML]{ECF4FF}\textbf{24.4\%} & \cellcolor[HTML]{ECF4FF}\textbf{31.5\%} & \cellcolor[HTML]{ECF4FF}\textbf{40.1\%} & \cellcolor[HTML]{ECF4FF}\textbf{22.2\%} & \cellcolor[HTML]{ECF4FF}\textbf{36.5\%} & \cellcolor[HTML]{ECF4FF}\textbf{43.1\%} & \cellcolor[HTML]{ECF4FF}\textbf{29.0\%} \\ \bottomrule
\end{tabular}
\end{table*}

\begin{itemize}

    \item \textbf{Code Coverage.} We measure three standard code coverage metrics: (1) \textit{Line Coverage}, which calculates the percentage of executed source code lines; (2) \textit{Function Coverage}, which measures the percentage of functions that are invoked at least once; and (3) \textit{Branch Coverage}, which tracks the percentage of conditional branches (e.g., if-else statements) that are executed. These metrics provide a holistic view of how thoroughly the generated test cases exercise the DBMS codebase.

    \item \textbf{Module-Level Coverage.} To understand the fine-grained impact of MIST on different DBMS components, we also measure coverage across four critical functional modules: (1) \textit{Parser}, which handles SQL parsing and syntax validation; (2) \textit{Optimizer}, which performs query optimization and plan generation; (3) \textit{Executor}, which executes the query plan and produces results; and (4) \textit{Storage}, which manages data storage and retrieval operations. This module-level analysis reveals which components benefit most from MIST's test generation strategy.

\end{itemize}

\subsection{Implementation Details}

All LLMs are downloaded from Hugging Face~\cite{Huggingface} and deployed locally using vLLM~\cite{kwon2023efficient} as the inference engine on 8 NVIDIA GPUs, with generation parameters configured as follows: maximum token length of 1,024, temperature of 0.2 for stable generation, and batch size of 1 with 3 concurrent workers to balance throughput and resource utilization. In Stage I, we generate 900 initial SQL test cases for each DBMS, with each test case incorporating 3 to 20 selected features from the hierarchical feature tree to ensure diversity. The range of 3 to 20 features is chosen through preliminary experiments to balance semantic richness and syntactic complexity: fewer than 3 features tend to produce overly simple queries with limited coverage, while more than 20 features often result in syntactically invalid SQL due to feature conflicts. In Stage II, we perform 600 mutation iterations to explore deeper execution paths and improve coverage incrementally. After generation, we apply post-processing to ensure the generated SQL can be correctly executed, including removing Markdown code block markers, stripping explanatory text and comments, extracting only valid SQL statements (DDL, DML, and DQL), and ensuring proper semicolon termination following standard practices in LLM-based code generation~\cite{joel2024survey,wang2023review,liu2024exploring}. For the MCTS mutation, we set the exploration weight $c$ in the UCT formula to 1.414 following the standard UCT formula recommendation~\cite{kocsis2006bandit}, and define an early-termination threshold of 50 new branches, which is empirically set to balance exploration efficiency and coverage improvement potential.

